\section{The cocharacter fixed scheme}\label{sec:cocharacters}

We now compare the fixed scheme for the chosen cocharacter with the fixed
scheme for the full topological torus.  This comparison uses only minuscule
generation, a root-support consequence of the loop-nilpotent Hall algebra,
and the regularity of the stability.  In particular, it does not use
Theorem~\ref{thm:full-torus} or assume coefficient surjectivity onto the
cocharacter fixed quotient.

If \(\nu\cdot\mathbf b\ne0\) for every
\(0<\mathbf b\leq\mathbf v\), Proposition~\ref{prop:minuscule} immediately
gives the comparison: every nontrivial positive or negative minuscule
generator has nonzero \(\nu\)-weight.  The argument below also covers
regular characters that lie on some of these dimension hyperplanes.

For example, take two disconnected vertices, no arrows, and
\(\mathbf v=\mathbf w=(1,1)\), with \(\nu=(1,-1)\).
Each Higgs factor is a point with free stable gauge action, so this
character is regular, although \(\nu\cdot(1,1)=0\).
The joint positive monopole is a product of monopoles of weights
\(1\) and \(-1\), and hence vanishes in the cocharacter fixed quotient:
weight zero alone does not imply survival.

To treat these characters, we adjoin a rank-one framing vertex.
If a nontrivial minuscule generator of the original Coulomb algebra
survived in its cocharacter fixed quotient, the corresponding class and
its complementary class would both survive in the enlarged cocharacter
fixed quotient.  Polynomial Hall-algebra generation would then decompose
both nonzero complementary dimension vectors into indecomposable dimension
vectors of weight zero for the extended character.
These two decompositions supply at least two blocks of a semistable direct
sum; after removing the one ineffective diagonal scalar, its stabilizer
remains positive-dimensional, contradicting regularity.

\subsection{Generation in indecomposable charges}

For a finite quiver \(\Gamma\), possibly with loops, write
\[
 \Delta_\Gamma
 =\{\,0\ne\boldsymbol\gamma\in\Nbb^{\Gamma_0}:
                  A_{\Gamma,\boldsymbol\gamma}(q)\ne0\,\}.
\]
Here \(A_{\Gamma,\boldsymbol\gamma}\) is the polynomial counting
absolutely indecomposable representations over finite fields introduced
in Kac's work \cite{Kac83}; polynomiality for quivers with self-loops
is explicitly recalled in \cite[Theorem~1(a) and Remark~4]{Kac19}.
We use \(\Delta_\Gamma\) only as the nonzero support of these polynomials,
without a loopless root-system interpretation.

The generation mechanism belongs to the cohomological Hall theory
initiated by Kontsevich--Soibelman \cite{KontsevichSoibelman11}.
Its BPS and PBW structure was developed by Davison--Meinhardt and
Davison \cite{DavisonMeinhardt20,Davison25AffineBPS}; the graded Kac
identity builds on Mozgovoy's generating-series formula and Davison's
cohomology of preprojective stacks
\cite{Mozgovoy11McKay,Davison23Preprojective}.
The difference-operator maps were developed in the work of
Finkelberg--Tsymbaliuk, Frassek--Tsymbaliuk, and Tsymbaliuk
\cite{FinkelbergTsymbaliuk19Integral,FrassekTsymbaliuk22,Tsymbaliuk23Difference}.
For arbitrary unframed quivers, Muthiah--Weekes relate Kac-polynomial
Hilbert series and shuffle algebras through limit quantized zastava
spaces \cite{MuthiahWeekes2026Zastavas}.
The zastava comparison provides related context.  Our inputs are the
polynomial-parameter loop-nilpotent package of Jindal--Negu\c{t}
\cite{JindalNegut26} and Tsymbaliuk's arbitrary-quiver operator maps
\cite[Theorem~5.7]{Tsymbaliuk23Difference}.
We first fix their product conventions and prove that the maps restrict
to polynomial coefficient rings.  This is needed before specializing
the independent parameters to the given flavor action.

Put
\[
 S=\Qbb[\hbar,u_a\mid a\in\Gamma_1],
 \qquad K=\operatorname{Frac}(S),
\]
with independent parameters.  Let \(\mathcal S_S^+\) be the polynomial
shuffle half of \cite[Definition~2.3]{JindalNegut26}, and denote its
product by \(*\).  It is graded by dimension vectors, and contains
\[
 e_{\mathbf n,g}
   =g(z_{ia})\prod_i\prod_{a,b=1}^{n_i}
                   (z_{ib}-z_{ia}+\hbar),
 \qquad
 g\in S[z_{ia}]^{\prod_iS_{n_i}}.
\]
For homogeneous elements of dimensions \(\mathbf n,\mathbf m\), use
the product
\begin{equation}\label{eq:hall-sign-product}
 E\star E'
   =(-1)^{\sum_{a:i\to j}n_jm_i}E*E'.
\end{equation}
Bilinearity of the exponent makes this product associative.
Let \(\mathcal S_S^-\) be its ordinary opposite algebra, with
\(f_{\mathbf n,g}\) denoting the same polynomial in that half.
A word changes under this twist by a sign depending only on the
dimensions of its factors; all homogeneous \(S\)-word spans are unchanged.

\begin{lemma}[Polynomial monopole maps]\label{lem:polynomial-monopole-maps}
For every unframed quiver \(\Gamma\) and dimension vector \(\mathbf d\),
the two halves above admit \(S\)-linear algebra maps to an algebra of
difference operators defined without inverting rotation or arrow parameters.
They send \(e_{\mathbf n,g}\) and \(f_{\mathbf n,g}\), for
\(\mathbf n\leq\mathbf d\), to nonzero constant multiples of the
positive and negative fundamental operators, respectively; they send
these elements to zero when \(\mathbf n\not\leq\mathbf d\).
After setting \(\hbar=0\), \(w_{ia}=x_{ia}\), and
\(u_a=-f_a\), where \(f_a\) is the prescribed flavor weight of arrow
\(a\), their images lie in \(A_{\Gamma,\mathbf d,F}\).
Dimension-vector degree becomes positive, respectively negative,
topological charge.  These maps preserve finite polynomial-coefficient
identities under this specialization.
\end{lemma}

\begin{proof}
We specify the operator algebra and the sign convention.
Let \(D_{ia}w_{jb}=(w_{jb}+\hbar\delta_{ij}\delta_{ab})D_{ia}\),
with all the \(D_{ia}\) commuting.  Allow coefficients in
\[
 S[w_{ia},(w_{ia}-w_{ib}+r\hbar)^{-1}:
                      a\ne b,\ r\in\Z].
\]
This ring is preserved by every shift \(D_{ia}^{\pm1}\).
For subsets \(A_i\subseteq\{1,\ldots,d_i\}\), write
\(w_A=(w_{ia})_{a\in A_i,i\in\Gamma_0}\).  Define
\begin{align*}
 E_{\mathbf n,g}
 &=\sum_{|A_i|=n_i}
   g(w_A)
   \frac{\displaystyle\prod_{a:i\to j}
             \prod_{\alpha\in A_i,\ \beta\notin A_j}
                (w_{i\alpha}-w_{j\beta}+u_a)}
        {\displaystyle\prod_i
             \prod_{\alpha\in A_i,\ \beta\notin A_i}
                (w_{i\alpha}-w_{i\beta})}
       \prod_{i,\alpha\in A_i}D_{i\alpha},\\
 F_{\mathbf n,g}
 &=\sum_{|A_i|=n_i}
   g(w_A-\hbar)
   \frac{\displaystyle\prod_{a:j\to i}
             \prod_{\alpha\in A_i,\ \beta\notin A_j}
                (w_{i\alpha}-w_{j\beta}-u_a)}
        {\displaystyle\prod_i
             \prod_{\alpha\in A_i,\ \beta\notin A_i}
                (w_{i\alpha}-w_{i\beta})}
       \prod_{i,\alpha\in A_i}D_{i\alpha}^{-1}.
\end{align*}
These are the fundamental formulas
\cite[Equations~(78)--(79)]{JindalNegut26}.
Let \(\mathcal C_S\) be the \(S\)-subalgebra generated by these operators
and the polynomials symmetric in \(w\) separately at each vertex.

We construct the maps first over \(K\), then prove that their polynomial
halves land in \(\mathcal C_S\).
Use Tsymbaliuk's maps on the full trigonometric shuffle algebra and its
opposite \cite[Theorem~5.7, Equations~(93) and~(96)]{Tsymbaliuk23Difference},
with no framing, and make the formal substitutions
\[
 q=e^{\varepsilon\hbar},\qquad
 t_a=e^{\varepsilon u_a},\qquad
 W_{ia}=e^{\varepsilon w_{ia}}.
\]
Here \(\varepsilon\) is an auxiliary formal variable, not \(\hbar\).
The multiplicative shift in the source satisfies
\(D^{\mathrm{trig}}_{ia}W_{ia}=q^{-1}W_{ia}D^{\mathrm{trig}}_{ia}\).
Thus its inverse becomes our positive shift \(D_{ia}\), while it
becomes \(D_{ia}^{-1}\) in the negative map.
The leading coefficient of the trigonometric shuffle kernel is
\[
 \left(\frac{x-\hbar}{x}\right)^{\delta_{ij}}
 \prod_{a:i\to j}[-(x+u_a)]
 \prod_{a:j\to i}(x+\hbar-u_a).
\]
It is precisely the kernel for \(\star\) in
Equation~\eqref{eq:hall-sign-product}: the incoming-arrow signs have
been removed from the kernel of \cite[Equation~(31)]{JindalNegut26}.
This distinction is necessary.  For the quiver consisting of one arrow
\(i\to j\), the untwisted product gives
\(e_{(1,0),1}*e_{(0,1),1}
   =-e_{(1,1),z_i-z_j+u_a}\), whereas the corresponding fundamental
operator product has a plus sign.

Here are the order calculations that justify taking leading coefficients,
including for loops and repeated gauge slots.
Put
\[
 \kappa(\mathbf n)=\sum_{a:i\to j}n_in_j,\qquad
 \ell_i=\#\{a:i\to i\},\qquad
 L_i^+=\sum_{a:i\to j}d_j-d_i+1,\qquad
 L_i^-=\sum_{a:j\to i}d_j-d_i+1.
\]
Lift a polynomial \(E\) of dimension \(\mathbf n\) to
\(\varepsilon^{\kappa(\mathbf n)}
 E((X_{ia}-1)/\varepsilon)\).
The cross-kernel order is
\(\kappa(\mathbf n+\mathbf m)-\kappa(\mathbf n)-\kappa(\mathbf m)\);
thus the leading source product is \(\star\).
In an operator summand, let \(m_{ia}\) be the length of the string
assigned to gauge slot \((i,a)\), so \(\sum_a m_{ia}=n_i\), and set
\(P_2=\sum_{i,a}\ell_i m_{ia}^2\).
The three prefactors in Tsymbaliuk's formulas have orders
\[
 \begin{array}{c|c}
 \text{factor}&\varepsilon\text{-order}\\ \hline
 \text{single-entry factors }Y\text{ or }Y'
     &L^\pm\cdot\mathbf n-\sum_i\ell_i n_i\\
 \text{inverse kernels within one string}
     &-P_2+\sum_i\ell_i n_i\\
 \text{factors between distinct slots}
     &-\kappa(\mathbf n)+P_2 .
 \end{array}
\]
Together with the lifted input, their total order is
\(L^\pm\cdot\mathbf n\), independently of the composition.
Multiply the operator image of the lift by
\(\varepsilon^{-L^\pm\cdot\mathbf n}\) and take its
\(\varepsilon^0\) coefficient to define the normalized leading map.
The remaining denominators have nonzero leading coefficients over
\(K(w)\).  There are only finitely many compositions.  Consequently
an input with zero normalized leading coefficient also has zero
normalized leading output.  The construction is independent of the lift.
Since \(L^\pm\cdot\mathbf n\) is additive, the normalized leading
maps preserve multiplication.  Multiplying them in dimension
\(\mathbf n\) by \(\hbar^{-|\mathbf n|}\) preserves that property.

The use of formal parameters in this argument requires no analytic
specialization theorem.  For real
\(\hbar<u_a<0\) and small positive \(\varepsilon\), the substitutions
satisfy Tsymbaliuk's parameter inequalities.
The monomial factors involving logarithms of torus parameters are
expanded near \(W=1\); they have leading coefficient one and rational
Taylor coefficients in \(\hbar,u_a,w_{ia}\).
Each coefficient of the multiplication identities is therefore a
rational identity valid on this nonempty real open set, hence valid
over \(K(w)\).  This justifies the formal leading maps, denoted
\(\Psi_+\) and \(\Psi_-\).

For the distinguished elements the maps are especially simple:
\begin{equation}\label{eq:corrected-monopole-images}
 \Psi_+(e_{\mathbf n,g})
   =(-1)^{L^+\cdot\mathbf n+\kappa(\mathbf n)}E_{\mathbf n,g},
 \qquad
 \Psi_-(f_{\mathbf n,g})=F_{\mathbf n,g}.
\end{equation}
Indeed, the same-color factor in \(e_{\mathbf n,g}\) vanishes on every
additive string of length at least two.  Only selected subsets of
distinct slots remain.  Its diagonal factors contribute
\(\hbar^{|\mathbf n|}\), which cancels the normalization
\(\hbar^{-|\mathbf n|}\); its off-diagonal factors cancel the
same-color denominators in the operator formula.  The factors for
arrows between two selected slots cancel as well.
In the positive formula the single-entry factors contribute
\((-1)^{\sum_i n_i(L_i^+-\ell_i)}\), and the internal-arrow
cancellations contribute
\((-1)^{\kappa(\mathbf n)-\sum_i\ell_i n_i}\), giving the stated sign.
In the opposite formula the selected variables are \(w_A-\hbar\);
the incoming-arrow factors cancel without these signs.
This proves both formulas, and proves that the image is zero when
\(\mathbf n\not\leq\mathbf d\).

By \cite[Theorem~2.4]{JindalNegut26}, the distinguished elements generate
\(\mathcal S_S^+\) over \(S\).  The twist preserves their word spans,
and the opposite statement follows by reversing words.
Equation~\eqref{eq:corrected-monopole-images} therefore proves that
both polynomial halves land in \(\mathcal C_S\).
Although we constructed the maps over \(K\), their restrictions are
thus defined over \(S\); no inverse of a parameter is needed in the
identities in \(\mathcal C_S\).

Finally specialize \(\hbar=0\), \(w_{ia}=x_{ia}\), and \(u_a=-f_a\).
Every allowed denominator becomes \(x_{ia}-x_{ib}\), with \(a\ne b\),
so this defines a homomorphism to
\(\operatorname{Frac}(P_F)[X_*(T)]\), sending each \(D_{ia}\)
to the monomial for its standard slot coweight.
For a selected-variable dressing \(g\), also denoting its specialization
by \(g\), direct comparison with Proposition~\ref{prop:minuscule} gives
\[
 M^+_{\mathbf n,g}
   =(-1)^{\sum_{a:i\to j}n_i(d_j-n_j)}
          E_{\mathbf n,g}\big|_{\mathrm{specialization}},
 \qquad
 M^-_{\mathbf n,g}
   =(-1)^{\sum_i n_i(d_i-n_i)}
          F_{\mathbf n,g}\big|_{\mathrm{specialization}}.
\]
The first sign comes from the outgoing matter factors; the second
comes from reversing the minuscule tangent weights.
The coefficient roots themselves are unchanged.
Thus every generator of \(\mathcal C_S\) specializes to an element
of \(A_{\Gamma,\mathbf d,F}\), and so does every image of the
polynomial halves.  Only this homomorphism is needed; no injectivity
or flatness after specialization is asserted.
\end{proof}

\begin{lemma}[Root-charge generation]\label{lem:root-charge-generation}
Let \(\Gamma\) be an unframed quiver with dimension vector \(\mathbf d\),
and let a flavor torus act by a character on each arrow.
Write \(A_{\Gamma,\mathbf d,F}\) for its commutative Coulomb algebra and
\(R_{\Gamma,\mathbf d,F}\) for its coefficient algebra.
For \(\epsilon\in\{+1,-1\}\), let
\(M_{\mathbf b,p}^{\epsilon}\) be a dressed minuscule class whose
coweight has \(b_i\) entries equal to \(\epsilon\) at vertex \(i\) and
all other entries zero, where \(0<\mathbf b\leq\mathbf d\).
There is a finite expression
\begin{equation}\label{eq:root-charge-words}
 M_{\mathbf b,p}^{\epsilon}
   =\sum_r c_r a_{r1}\cdots a_{rk_r},
 \qquad c_r\in R_{\Gamma,\mathbf d,F},
\end{equation}
in which every factor is topologically homogeneous and
\[
 \deg_{\mathrm{top}}a_{rh}=\epsilon\boldsymbol\gamma_{rh},
 \qquad
 \boldsymbol\gamma_{rh}\in\Delta_\Gamma,
 \qquad
 \sum_{h=1}^{k_r}\boldsymbol\gamma_{rh}=\mathbf b.
\]
The factors need not themselves be minuscule classes.
\end{lemma}

\begin{proof}
\emph{Polynomial Hall-algebra input.}
Use the independent polynomial ring \(S\) defined above.  This is
the equivariant coefficient ring $H^*(BT_{\mathrm{univ}};\Qbb)$ for
$T_{\mathrm{univ}}=(\C^\times)^{1+|\Gamma_1|}$, with independent
rotation and arrow characters.  It lies in the geometric parameter
setting of \cite[Equation~(30) and Remark~2.6]{JindalNegut26}; no
extension to arbitrary field-valued exponential parameters is needed.
On the
tripled quiver use weights \(-u_a\), \(u_a-\hbar\), and \(\hbar\)
on an original arrow, its reverse, and an added loop, respectively.
These choices satisfy all the genericity hypotheses of
\cite[\S2.2]{JindalNegut26}; in particular there is no nonzero
nonnegative integral relation among \(u_a,\hbar-u_a\), and
\(u_a\notin\Z\hbar\) for every original loop.

We use the following polynomial-form results of that source:
the loop-nilpotent Hall algebra identifies with its shuffle half
\cite[Proposition~3.7]{JindalNegut26};
its BPS descendants generate over \(S\)
\cite[Theorem~3.10, Equation~(109), and
Corollary~3.15]{JindalNegut26};
and the BPS comparison and Kac formula show that the nonzero
dimension-vector support of these generators is \(\Delta_\Gamma\)
\cite[Proposition~3.14 and \S3.10]{JindalNegut26}.
Here the Kac identity retains the full cohomological grading: the
coefficient of $t^{d/2}$ is $(-1)^d\dim H^d$.
A zero Kac polynomial therefore forces every graded piece of the
corresponding BPS space to vanish.  Proposition~3.14 of that source
transfers this vanishing to the loop-nilpotent BPS space.
The concrete twist in Equation~\eqref{eq:hall-sign-product} preserves
these homogeneous word spans.
Here polynomial form refers to equivariant parameters over \(\Qbb\),
not integral cohomology.

Apply the generation statement to \(e_{\mathbf b,g}\) or
\(f_{\mathbf b,g}\), and retain dimension-vector degree \(\mathbf b\).
All nonconstant Hall degrees are nonnegative, so the result is a
finite \(S\)-linear combination of products of BPS descendants with
degrees in \(\Delta_\Gamma\) summing to \(\mathbf b\).
Applying Lemma~\ref{lem:polynomial-monopole-maps} gives the corresponding
finite identity in \(\mathcal C_S\).
Its coefficients lie in \(S\); no parameter from \(S\) is inverted
in this identity.
Every factor has degree at most \(\mathbf b\leq\mathbf d\); no assertion
about vanishing of entire Hall components outside this bound is needed.

\emph{Specialization.}
Now extend scalars to \(\C\) and apply the specialization homomorphism
of Lemma~\ref{lem:polynomial-monopole-maps}.
Each finite identity becomes an identity in \(A_{\Gamma,\mathbf d,F}\).
The signs in the images of the distinguished elements are nonzero
constants and can be absorbed in its coefficients.
This operation is performed on the abstract polynomial-form algebra and
its map, not on the ambient polynomial representation of the shuffle
algebra.  An element whose ambient representative is divisible by
\(\hbar\) need not belong to \(\hbar\) times the polynomial shuffle
half, and is therefore not discarded before applying the map.

\emph{General dressings.}
It remains to allow an arbitrary stabilizer-invariant dressing \(p\).
At each vertex, the polynomial algebra symmetric separately in the
selected and complementary variables is generated by the polynomials
symmetric in all variables together with those symmetric in the
selected variables.  Indeed, each complementary power sum is the total
power sum minus the selected power sum.  Thus \(p\) is a finite sum
of coefficient polynomials times selected-variable dressings \(g\).
The already proved identities for \(g\) give
Equation~\eqref{eq:root-charge-words}.
\end{proof}

\subsection{Indecomposables and central moment values}

We record the geometric consequence of the charges in
Lemma~\ref{lem:root-charge-generation}.  For a dimension vector
\(\boldsymbol\gamma\), put
\[
 G_{\boldsymbol\gamma}=\prod_{i\in\Gamma_0}\GL_{\gamma_i},
 \qquad
 N_{\boldsymbol\gamma}=\bigoplus_{a:i\to j}
                     \Hom(\C^{\gamma_i},\C^{\gamma_j}).
\]
For a real vector \(\boldsymbol\theta=(\theta_i)\), let
\(\tau_{\boldsymbol\theta}\in\Lie(G_{\boldsymbol\gamma})^*\)
be the central functional
\[
 \tau_{\boldsymbol\theta}(\xi)
       =\sum_i\theta_i\operatorname{tr}(\xi_i).
\]
Choose standard Hermitian forms, and write
\((\mu_{\mathbb R},\mu_{\mathbb C})\) for the flat hyperkähler moment
maps on \(T^*N_{\boldsymbol\gamma}\).  We normalize the real parameter
so that an integral character \(\boldsymbol\theta\) corresponds to
real central level \(\tau_{\boldsymbol\theta}\) in the affine
Kempf--Ness correspondence.  With any other common normalization,
the same fixed scalar multiplies every level below.

\begin{lemma}\label{lem:root-moment}
If \(\boldsymbol\gamma\in\Delta_\Gamma\) and
\(\boldsymbol\theta\cdot\boldsymbol\gamma=0\), then
\[
 \mu_{\mathbb C}^{-1}(0)\cap
 \mu_{\mathbb R}^{-1}(\tau_{\boldsymbol\theta})\ne\varnothing.
\]
The scalar circle acting simultaneously on all nonzero vertex spaces
fixes every point of this representation.
\end{lemma}

\begin{proof}
First, \(\Gamma\) has an indecomposable complex representation \(x\)
of dimension \(\boldsymbol\gamma\).
For completeness, the required passage from Kac support to characteristic
zero can be made by constructibility.  A representation is decomposable
over an algebraic closure exactly when its endomorphism algebra contains
an idempotent other than \(0\) and \(1\).  The equations for such
idempotents, together with these two open exclusions, define a
constructible incidence condition over \(\Z\).
Chevalley's theorem therefore gives a constructible locus of
geometrically indecomposable representations.  A nonzero Kac polynomial
gives points of this locus in infinitely many residue characteristics.
Its generic fiber cannot be empty, since a constructible subset of
\(\Spec\Z\) that omits the generic point is finite.  Base change of
the nonempty characteristic-zero locus to \(\C\) gives \(x\).

The following trace-annihilator argument is the lifting criterion of
Crawley--Boevey \cite[Lemma~3.1 and Theorem~3.3]{CrawleyBoevey01}; we
recall it before passing to the real moment level.
Every endomorphism \(e\) of \(x\) has the form \(c\id+n\), where \(n\)
is nilpotent.  Indeed, its generalized eigenspaces are subrepresentations,
so indecomposability permits only one eigenvalue.
Consequently
\[
 \tau_{\boldsymbol\theta}(e)
 =c\,\boldsymbol\theta\cdot\boldsymbol\gamma=0.
\]
The cotangent moment map with \(x\) fixed is the transpose of the
infinitesimal orbit map:
\[
 N_{\boldsymbol\gamma}^*\longrightarrow
       \Lie(G_{\boldsymbol\gamma})^*,
 \qquad
 y\longmapsto\bigl(\xi\longmapsto y(\xi x)\bigr).
\]
Its image is the annihilator of
\(\End_\Gamma(x)=\ker(\xi\mapsto\xi x)\).  Hence the complex moment
fiber at \(\tau_{\boldsymbol\theta}\) is nonempty, as is the fiber at
any common scalar multiple of that value.

This fiber is a closed invariant affine subset, since the moment value
is central.  The closure of any orbit in it contains an orbit closed
in the ambient affine representation.  By the zero-character affine
Kempf--Ness theorem, that orbit meets \(\mu_{\mathbb R}^{-1}(0)\);
see \cite[Theorem~6.1 and Corollary~6.2]{King94}.
Thus there is a point with hyperkähler moment value
\((0,\tau_{\boldsymbol\theta},0)\), using the real and imaginary
parts of the complex moment map and adjusting the common scalar to
the chosen normalization.

The quaternionic rotation on the flat cotangent representation commutes
with the compact gauge group and rotates its three moment maps.
A single rotation sends this value to
\((\tau_{\boldsymbol\theta},0,0)\)
\cite[\S3(D)]{HKLR87}.
This proves the asserted nonemptiness.  The simultaneous scalar acts
trivially on all arrow spaces and their duals, proving the last statement.
\end{proof}

\subsection{The framing vertex and complementary monopoles}

\begin{proposition}\label{prop:cocharacter-fixed}
Suppose that every \(\nu\)-semistable point of \(\mu^{-1}(0)\) is stable
and has trivial \(G\)-stabilizer.  Then the natural quotient map
\[
 A_F/\langle (A_F)_m:m\ne0\rangle
       \longrightarrow B_F,
\]
where \(m\) denotes \(\nu\)-weight, is an isomorphism.
This is an equality of quotients of \(A_F\), preserving the coefficient
map and both gradings.
\end{proposition}

\begin{proof}
Lemma~\ref{lem:trivial-summands} allows us to suppose that \(F\) scales
the individual arrows and framing rows.  Adding or removing the
gauge-trivial cotangent summands in that lemma preserves both regularity
and all the fixed quotients under consideration.
Write
\[
 \nu((g_i))=\prod_i\det(g_i)^{\nu_i}.
\]
Adjoin one vertex \(\infty\) of dimension one and \(w_i\) arrows
\(i\to\infty\), keeping the original arrows.  Denote this unframed
quiver by \(\widehat Q\) and its dimension vector by
\(\mathbf d=(\mathbf v,1)\).
Its representation is the original \(N\), after a choice of basis
in each framing space.

There is a group isomorphism
\[
 G\times\C^\times\longrightarrow G_{\mathbf d},
 \qquad ((g_i),c)\longmapsto((cg_i),c).
\]
The second factor acts trivially on \(N\).  The BFN space of triples
and its convolution therefore split, giving
\begin{equation}\label{eq:central-splitting}
 A_{\widehat Q,\mathbf d,F}
       \cong A_F\otimes_\C\C[\xi,z,z^{-1}].
\end{equation}
Here \(\xi\) is the coefficient parameter of the extra scalar and
\(z\) is its fundamental Laurent monopole, of cohomological degree zero.
This \(z\) is unrelated to the degree-two added-loop mass in
Section~\ref{sec:deformations}.  In the original vertex
coordinates, \(z\) has topological charge \(\mathbf d\): it translates
every gauge coweight entry, including the entry at \(\infty\), by one.

Extend the character by
\[
 \widehat\nu_i=\nu_i\quad(i\in Q_0),
 \qquad
 \widehat\nu_\infty=-\sum_i\nu_i v_i.
\]
Then \(\widehat\nu\cdot\mathbf d=0\), and under the displayed group
isomorphism \(\widehat\nu\) is \(\nu\) on \(G\) and trivial on the
extra scalar.  In particular, both \(\xi\) and \(z^{\pm1}\) have
\(\widehat\nu\)-weight zero.  Taking the cocharacter fixed quotient
in Equation~\eqref{eq:central-splitting} consequently gives
\[
 \bigl(A_{\widehat Q,\mathbf d,F}\bigr)_{\mathrm{fix},\widehat\nu}
 \cong
 \bigl(A_F\bigr)_{\mathrm{fix},\nu}
                 \otimes_\C\C[\xi,z,z^{-1}],
\]
where the subscript denotes the quotient by nonzero cocharacter
weights.  This extension is faithfully flat.

Let \(M_{\mathbf b',p}^{+}\) be a nontrivial positive minuscule generator
of \(A_F\), and put \(\mathbf b=(\mathbf b',0)\).
If \(\widehat x_{ia}\) denotes an enlarged gauge Chern root, the group
isomorphism above identifies the original root with
\(x_{ia}=\widehat x_{ia}-\xi\).
Let \(p^{\uparrow}\) be the polynomial obtained from \(p\) by these
substitutions, leaving every flavor parameter unchanged.
The original generator is then the positive minuscule class
\(M_{\mathbf b,p^{\uparrow}}^{+}\).
Central translation also expresses it using the complementary negative
coweight:
\begin{equation}\label{eq:complementary-monopoles}
 M_{\mathbf b,p^{\uparrow}}^{+}
       =z\,M_{\mathbf d-\mathbf b,p^{\uparrow}}^{-}.
\end{equation}
The two coweights have the same Weyl stabilizer, so the same polynomial
\(p^{\uparrow}\) is an admissible dressing on both sides.
It need not depend only on the variables selected by the complementary
negative coweight; Lemma~\ref{lem:root-charge-generation} allows all
stabilizer-invariant dressings.
The equality follows directly from
Equation~\eqref{eq:minuscule}: all matter weights pair trivially
with the central coweight, and central translation identifies the
minuscule orbits and their equivariant tangent bundles.
Equivalently it follows from convolution with the invertible central
class.  Thus it holds with compatible fundamental-class conventions;
rescaling minuscule generators changes only a nonzero constant.

Suppose that this generator has nonzero image in the cocharacter
fixed quotient.  Apply Lemma~\ref{lem:root-charge-generation} in the
enlarged quiver to \(M_{\mathbf b,p^{\uparrow}}^{+}\).
Any word containing a factor of nonzero \(\widehat\nu\)-weight
vanishes in that quotient.  A nonzero image therefore requires a
decomposition
\[
 \mathbf b=\sum_a\boldsymbol\alpha_a,
 \qquad
 \boldsymbol\alpha_a\in\Delta_{\widehat Q},
 \qquad
 \widehat\nu\cdot\boldsymbol\alpha_a=0.
\]
Since \(z\) is an invertible element of weight zero,
Equation~\eqref{eq:complementary-monopoles} implies that
\(M_{\mathbf d-\mathbf b,p^{\uparrow}}^{-}\) also survives.  Applying the negative
half of the same lemma yields
\[
 \mathbf d-\mathbf b=\sum_b\boldsymbol\beta_b,
 \qquad
 \boldsymbol\beta_b\in\Delta_{\widehat Q},
 \qquad
 \widehat\nu\cdot\boldsymbol\beta_b=0.
\]
Both sums are nonempty: \(\mathbf b\ne0\), while its \(\infty\)-entry
is zero and that of \(\mathbf d\) is one.

Together these decompositions express \(\mathbf d\) as a sum of at
least two nonzero dimension vectors
\(\boldsymbol\gamma_1,\ldots,\boldsymbol\gamma_k\) in
\(\Delta_{\widehat Q}\), each orthogonal to \(\widehat\nu\).
For each summand choose the point furnished by
Lemma~\ref{lem:root-moment}, with real central parameter
\(\widehat\nu\).  Use orthogonal decompositions of the vertex spaces
and take the direct sum of these points.  Moment maps are block
diagonal on direct sums, so the resulting point has
\[
 \mu_{\mathbb C}=0,\qquad
 \mu_{\mathbb R}=\tau_{\widehat\nu}.
\]
Its stabilizer contains \(k\) independent block scalar circles.

The enlarged reduction is exactly the original framed reduction:
the extra diagonal scalar is ineffective, its moment equation is
redundant, and \(\widehat\nu\) is trivial on it.
After dividing by this scalar, the stabilizer still contains a torus
of dimension \(k-1>0\).
By the affine Kempf--Ness correspondence at the character level, the
resulting point is \(\nu\)-semistable in the original zero complex
moment fiber \cite[Theorem~6.1 and Corollary~6.2]{King94}.
This contradicts the assumed freeness of every such point.
The argument also applies when the stable locus is empty: it would
then produce a semistable point where none exists.

Every nontrivial positive minuscule generator therefore vanishes.
For a negative generator use the identity
\(M_{\mathbf b,p^{\uparrow}}^{-}
  =z^{-1}M_{\mathbf d-\mathbf b,p^{\uparrow}}^{+}\)
and interchange the two halves in the argument; it vanishes as well.
By Proposition~\ref{prop:minuscule}, the ideal they generate is the ideal
generated by the nonzero full topological-charge spaces.  It is therefore
contained in the ideal generated by the nonzero \(\nu\)-weight spaces.
For \(m\ne0\), the \(\nu\)-weight space \((A_F)_m\) is a sum of nonzero
topological-charge spaces.  This gives the opposite inclusion.
The ideals coincide, proving the proposition.
If all gauge ranks are zero there are no such generators and the
same conclusion is immediate.
\end{proof}
